\documentclass[11pt]{article}

\usepackage{fullpage}
\usepackage{amsmath, amssymb, bm, cite, epsfig, psfrag}
\usepackage{graphicx}
\usepackage{float}
\usepackage{amsthm}
\usepackage{amsfonts}
\usepackage{cite}
\usepackage{hyperref}
\usepackage{pgf,tikz}
\usepackage{enumitem}
\usepackage{mathtools}
\usepackage{siunitx}
\usepackage{../../styles/course}


\begin{document}

\title{Introduction to Hardware Design\\
Arrays:  Figures and Code Snippets}
\author{Profs. Sundeep Rangan, Siddharth Garg}
\date{}

\maketitle
\subsection{array}

\begin{tikzpicture}
    \def\nrow{2}
    \def\ncol{3}

    \pgfmathtruncatemacro{\lastrow}{\nrow-1}
    \pgfmathtruncatemacro{\lastcol}{\ncol-1}

    \tikzset{
        multbox/.style={
            draw,
            rectangle,
            fill=blue!20,
            minimum width=1.8cm,
            minimum height=1.0cm
        }
    }
    \tikzset{
        xbox/.style={
            draw,
            rectangle,
            fill=green!20,
            minimum width=1.8cm,
            minimum height=1.0cm
        }
    }
    \tikzset{
        ybox/.style={
            draw,
            rectangle,
            fill=purple!20,
            minimum width=1.8cm,
            minimum height=1.0cm
        }
    }

    \node[xbox] (X0) {$X_{n,0}$};
    \foreach \i in {1,...,\lastrow} {
        \pgfmathtruncatemacro{\prevrow}{\i-1}
        \node[xbox,below=1cm of X\prevrow] (X\i) {$X_{n,\i}$};
    }

    \foreach \i in {0,...,\lastrow} {
        \node[multbox, right=1cm of X\i] (W\i0) {$W_{\i,0}$};
        \foreach \j in {1,...,\lastcol} {
            \pgfmathtruncatemacro{\prevcol}{\j-1}
            \node[multbox, right=1cm of W\i\prevcol] (W\i\j) {$W_{\i,\j}$};
        }

        \foreach \j in {0,...,\lastcol} {
            \pgfmathsetmacro{\arcbend}{1.1 + 0.1*\j}
            \draw [thick, green!50!black, ->] (X\i.north) to[out=45,in=135,looseness=0.3] (W\i\j.north);
        }
    }

    \foreach \j in {0,...,\lastcol} {
        \node[ybox, below=1cm of W\lastrow\j] (Y\j) {$Y_{n,\j}$};
        \foreach \i in {0,...,\lastrow} {
            \draw [thick, purple, ->] (W\i\j.east) to[out=--45,in=45,looseness=0.3] (Y\j.east);
        }
    }

\end{tikzpicture}


\subsection{Fanout}

Consider simple Vitis code:

\begin{ccode}
for (int i = 0; i < N; i++) {
#pragma HLS PIPELINE II=1
    for (int j = 0; j < M; j++) {
#pragma HLS UNROLL
        y[i][j] = f(x[i][j], b);
    }
}
\end{ccode}

Tis can be visualized as:

\pagebreak


\begin{figure}
\begin{tikzpicture}
    \def\nrow{4}

    \pgfmathtruncatemacro{\lastrow}{\nrow-1}
   
    \tikzset{
        bbox/.style={
            draw,
            rectangle,
            fill=blue!20,
            minimum width=1.8cm,
            minimum height=1.0cm
        }
    }
    \tikzset{
        fbox/.style={
            draw,
            rectangle,
            fill=green!20,
            minimum width=1.8cm,
            minimum height=1.0cm
        }
    }

    \node[bbox] (b) {$b$};

    \foreach \i in {0,...,\lastrow} {
        \ifnum\i>0
            \pgfmathtruncatemacro{\prevrow}{\i-1}
            \node[fbox,right=1cm of f\prevrow] (f\i) {$f(x_{n\i},b)$};
        \else
            \node[fbox,below=1cm of b,xshift=-1cm] (f\i) {$f(x_{n\i},b)$};
        \fi
        \node[below=0.7cm of f\i,xshift=-0.3cm] (x\i) {$x_{n\i}$};
        \node[below=1.3cm of f\i,xshift=0.3cm] (y\i) {$y_{n\i}$};
        \draw [->] (x\i.north) -- (x\i.north |- f\i.south);
        \draw [->] (y\i.north |- f\i.south)-- (y\i.north);
        \draw [->,thick,green!50!black] (b) -- (f\i.north);
        }   
\end{tikzpicture}
\end{figure}


\begin{ccode}
bcopy[0] = b;
for (int j=1; j < M; j++) {
    bcopy[j] = bcopy[j-1];
}

for (int i = 0; i < N; i++) {
#pragma HLS PIPELINE II=1
    for (int j = 0; j < M; j++) {
#pragma HLS UNROLL
        y[i][j] = f(x[i][j], bcopy[j]);
    }
}
\end{ccode}

\begin{tikzpicture}
    \def\nrow{4}

    \pgfmathtruncatemacro{\lastrow}{\nrow-1}
   
    \tikzset{
        bbox/.style={
            draw,
            rectangle,
            fill=blue!20,
            minimum width=1.8cm,
            minimum height=1.0cm
        }
    }
    \tikzset{
        fbox/.style={
            draw,
            rectangle,
            fill=green!20,
            minimum width=1.8cm,
            minimum height=1.0cm
        }
    }

    \node[bbox] (b) {$b$};

    \node[bbox,right=1cm of b] (b0) {$b^{\subsf copy}_{m}$};
    \draw [->] (b.east) -- (b0.west);
    \foreach \i in {1,...,\lastrow} {
        \pgfmathtruncatemacro{\prevrow}{\i-1}
        \node[bbox,right=1cm of b\prevrow] (b\i) {$b^{\subsf copy}_{\i}$};
        \draw [->] (b\prevrow.east) -- (b\i.west);
    }

    \foreach \i in {0,...,\lastrow} {
        \node[fbox,below=1cm of b\i] (f\i) {$f(x_{n\i},b^{\subsf copy}_{\i})$};
        \node[below=0.7cm of f\i,xshift=-0.3cm] (x\i) {$x_{n\i}$};
        \node[below=1.3cm of f\i,xshift=0.3cm] (y\i) {$y_{n\i}$};
        \draw [->] (x\i.north) -- (x\i.north |- f\i.south);
        \draw [->] (y\i.north |- f\i.south)-- (y\i.north);
        \draw [->] (b\i.south) -- (f\i.north);
        }   
\end{tikzpicture}


\end{document}






